% sections/questions.tex -- research questions, provisional argument and terms (shared).
\textbf{Central question.} How do GenAI systems reorganize interpretive
authority and collective authorship in collaborative artistic practice?

\begin{enumerate}[label=\textbf{SQ\arabic*.}, leftmargin=*]
  \item \emph{Translation and cultural legibility.} How do selected models,
        interfaces, defaults, moderation systems and facilitation
        arrangements privilege, amplify, flatten, standardize, mistranslate,
        misrecognize, refuse or pass over participants' self-identified and
        culturally situated contributions? (Studio records; episodes;
        transformation histories; platform documentation.)
  \item \emph{Negotiation and governance.} How do collaborators negotiate
        uptake, revision, refusal, attribution, credit, responsibility and
        collective decisions among contributors, the group, the
        artist-facilitator, models, platforms and institutions? (Field
        notes; group rules; version histories; focus groups.)
  \item \emph{Research-creation and public form.} How can collaborative
        artistic methods and the form of a public artwork materialize
        otherwise hidden distributions of interpretive authority and leave
        them open to debate and revision? (Scores; versions of the work;
        studio journal.)
\end{enumerate}

The exhibition is a creative, analytical and dissemination outcome. Because
audience data will not be collected, the study will not make empirical
claims about audience reception.

\textbf{Provisional argument.} GenAI does not merely add a nonhuman
collaborator to collective art-making; it may reorganize the conditions
under which contributions become visible, legible, transformable,
contestable, attributable, creditable, ownable and accountable. The research
may confirm, complicate or reject this proposition, and it treats refusal,
silence, failure, non-use, withdrawal, disagreement and unresolved parallel
versions as meaningful outcomes.

\textbf{Terms.} \emph{Participation} means taking part in consequential
decisions, not merely receiving access to or interacting with a system.
\emph{Recognition} concerns whether a contribution remains perceptible and
attributable within the shared work. \emph{Interpretive authority} is the
capacity to influence or decide what a contribution means, how it is
transformed and what it becomes. \emph{Attribution} or \emph{credit} is the
public acknowledgement of a contribution. \emph{Authorship} is a socially
and institutionally negotiated claim that a person or group substantively
shaped the work. \emph{Ownership} is a legal or contractual claim.
\emph{Responsibility} concerns who is expected to answer for the work and
its consequences. \emph{Technical provenance} is a system's record of
inputs, transformations or versions: evidence, but not a definitive record
of social authorship, cultural recognition or responsibility.
\emph{Culturally situated contributions} include languages, references,
idioms, memories, aesthetic conventions, performance practices and forms of
community knowledge that participants themselves identify as significant.
The project will not judge cultural authenticity or treat anyone as
representing a culture or community.
